\documentclass[12pt,a4paper]{article}

% ---- Encoding & Font ----
\usepackage{iftex}
\ifPDFTeX
  \usepackage[utf8]{inputenc}
  \usepackage[T1]{fontenc}
\else
  \usepackage{fontspec}  % XeTeX/LuaTeX (tectonic): Unicode nativo (·, —, ¿, ª, §)
\fi
\usepackage[spanish]{babel}
\usepackage{csquotes}

% ---- Page layout ----
\usepackage[top=2.5cm, bottom=2.5cm, left=2.5cm, right=2.5cm]{geometry}
\usepackage{setspace}
\onehalfspacing

% ---- Math ----
\usepackage{amsmath, amssymb, amsthm}
\usepackage{mathtools}
\usepackage{physics}
\usepackage{esint}  % \oiint

% ---- Graphics ----
\usepackage{graphicx}
\usepackage{tikz}
\usepackage{float}
\usepackage{caption}

% ---- Tables ----
\usepackage{array}
\usepackage{booktabs}
\usepackage{tabularx}
\usepackage{colortbl}

% ---- Colours ----
\usepackage{xcolor}
\definecolor{notebg}{HTML}{FFF3CD}
\definecolor{noteborder}{HTML}{FFC107}
\definecolor{boxred}{HTML}{F8D7DA}
\definecolor{boxredborder}{HTML}{F5C6CB}

% ---- Boxes ----
\usepackage{mdframed}
\usepackage{tcolorbox}
\tcbuselibrary{most}

\newtcolorbox{keybox}[1][]{
  colback=blue!5,
  colframe=blue!70!black,
  arc=4pt,
  boxrule=1pt,
  fonttitle=\bfseries,
  title={#1}
}

\newtcolorbox{warnbox}[1][]{
  colback=boxred,
  colframe=boxredborder,
  coltitle=black!75,
  arc=4pt,
  boxrule=1pt,
  fonttitle=\bfseries,
  title={#1}
}

\newtcolorbox{notebox}[1][]{
  colback=notebg,
  colframe=noteborder,
  coltitle=black!75,
  arc=4pt,
  boxrule=1pt,
  fonttitle=\bfseries,
  title={#1}
}

% ---- Hyperlinks & TOC ----
\usepackage[hidelinks]{hyperref}
\usepackage{bookmark}

% ---- Enumerate ----
\usepackage{enumitem}

% ---- Miscellaneous ----
\usepackage{icomma}
\usepackage{siunitx}
\usepackage{textcomp}
\usepackage[bottom]{footmisc}

% ---- Diagramas (circuitos, gráficas y esquemas) ----
\usepackage[europeanresistors]{circuitikz}
\usepackage{pgfplots}
\pgfplotsset{compat=1.18}
\usetikzlibrary{babel}  % babel español activa `>`; esto evita romper las flechas `->`
% courses/ElecMag2024/Clases/assets/tikzstyles.tex
% ---------------------------------------------------------------------------
% Estilos compartidos para los diagramas del curso Electromagnetismo (2024).
% Fuente única del "look" visual (colores, grosores, escalas) para que todas
% las figuras (TikZ / pgfplots / CircuiTikz) sean consistentes entre clases.
%
% Es una COPIA de courses/Fisica3-2015/Clases/assets/tikzstyles.tex: mismo
% dominio, misma paleta. Copia y no symlink para que cada curso sea
% autocontenido y la paleta de Física III pueda evolucionar aparte.
%
% Uso: la clase debe cargar antes `tikz` y `pgfplots` (y `circuitikz` si la
%      clase tiene circuitos), y luego % courses/ElecMag2024/Clases/assets/tikzstyles.tex
% ---------------------------------------------------------------------------
% Estilos compartidos para los diagramas del curso Electromagnetismo (2024).
% Fuente única del "look" visual (colores, grosores, escalas) para que todas
% las figuras (TikZ / pgfplots / CircuiTikz) sean consistentes entre clases.
%
% Es una COPIA de courses/Fisica3-2015/Clases/assets/tikzstyles.tex: mismo
% dominio, misma paleta. Copia y no symlink para que cada curso sea
% autocontenido y la paleta de Física III pueda evolucionar aparte.
%
% Uso: la clase debe cargar antes `tikz` y `pgfplots` (y `circuitikz` si la
%      clase tiene circuitos), y luego % courses/ElecMag2024/Clases/assets/tikzstyles.tex
% ---------------------------------------------------------------------------
% Estilos compartidos para los diagramas del curso Electromagnetismo (2024).
% Fuente única del "look" visual (colores, grosores, escalas) para que todas
% las figuras (TikZ / pgfplots / CircuiTikz) sean consistentes entre clases.
%
% Es una COPIA de courses/Fisica3-2015/Clases/assets/tikzstyles.tex: mismo
% dominio, misma paleta. Copia y no symlink para que cada curso sea
% autocontenido y la paleta de Física III pueda evolucionar aparte.
%
% Uso: la clase debe cargar antes `tikz` y `pgfplots` (y `circuitikz` si la
%      clase tiene circuitos), y luego \input{../assets/tikzstyles.tex}
% (compilar desde el directorio de la clase para que la ruta relativa resuelva).
% ---------------------------------------------------------------------------

% ---- Paleta (alineada con las cajas del preámbulo: keybox/notebox/warnbox) ----
\definecolor{figblue}{HTML}{1F4E79}   % azul principal (líneas, keybox)
\definecolor{figred}{HTML}{C0392B}    % rojo de énfasis (warnbox)
\definecolor{figamber}{HTML}{B8860B}  % ámbar (notebox)
\definecolor{figgray}{HTML}{555555}   % auxiliares, ejes, guías

% ---- CircuiTikz: escala y grosor de los componentes ----
% Guarda \ifdefined: la electrostática (Clases 1-14) no carga `circuitikz`, y
% sin la guarda el \input revienta con `Undefined control sequence \ctikzset`.
% Cuando el curso llegue a circuitos, el bloque se activa solo.
\ifdefined\ctikzset
  \ctikzset{
    bipoles/thickness=1,
    resistors/scale=0.9,
    capacitors/scale=0.9,
  }
\fi

% ---- TikZ: estilos reutilizables ----
% Nota: las puntas se declaran como `-{Latex[...]}` y no como `->`. La forma
% con llaves NO contiene un `>` literal, así que es inmune al `>` activo que
% introduce babel español (ver CLAUDE.md §7.3.3).
\usetikzlibrary{arrows.meta,calc,angles,quotes}

\tikzset{
  figline/.style={figblue, line width=0.9pt},
  figaux/.style={figgray, dashed, line width=0.6pt},
  % vectores
  figvec/.style={figblue, line width=0.9pt, -{Latex[length=2.2mm,width=1.6mm]}},
  figvecres/.style={figred, line width=1.3pt, -{Latex[length=2.6mm,width=1.9mm]}},
  figvecaux/.style={figgray, line width=0.7pt, -{Latex[length=1.8mm,width=1.3mm]}},
  figaxis/.style={figgray, line width=0.6pt, -{Latex[length=2mm,width=1.4mm]}},
  % cargas puntuales
  figcarga/.style={circle, inner sep=0pt, minimum size=5.4pt, draw=none},
  figpos/.style={figcarga, fill=figred},
  figneg/.style={figcarga, fill=figblue},
  figneutra/.style={figcarga, fill=figgray},
  % etiquetas
  figlbl/.style={font=\small},
  figlblaux/.style={font=\footnotesize, figgray},
}

% ---- pgfplots: estilo base de las gráficas ----
\pgfplotsset{
  emfig/.style={
    width=11cm, height=6.5cm,
    axis lines=left,
    xlabel style={font=\small},
    ylabel style={font=\small},
    tick label style={font=\footnotesize},
    every axis plot/.append style={line width=1pt},
    grid=major,
    grid style={figgray!20},
    legend cell align=left,
    legend style={font=\footnotesize, draw=figgray!40},
  },
}

% (compilar desde el directorio de la clase para que la ruta relativa resuelva).
% ---------------------------------------------------------------------------

% ---- Paleta (alineada con las cajas del preámbulo: keybox/notebox/warnbox) ----
\definecolor{figblue}{HTML}{1F4E79}   % azul principal (líneas, keybox)
\definecolor{figred}{HTML}{C0392B}    % rojo de énfasis (warnbox)
\definecolor{figamber}{HTML}{B8860B}  % ámbar (notebox)
\definecolor{figgray}{HTML}{555555}   % auxiliares, ejes, guías

% ---- CircuiTikz: escala y grosor de los componentes ----
% Guarda \ifdefined: la electrostática (Clases 1-14) no carga `circuitikz`, y
% sin la guarda el \input revienta con `Undefined control sequence \ctikzset`.
% Cuando el curso llegue a circuitos, el bloque se activa solo.
\ifdefined\ctikzset
  \ctikzset{
    bipoles/thickness=1,
    resistors/scale=0.9,
    capacitors/scale=0.9,
  }
\fi

% ---- TikZ: estilos reutilizables ----
% Nota: las puntas se declaran como `-{Latex[...]}` y no como `->`. La forma
% con llaves NO contiene un `>` literal, así que es inmune al `>` activo que
% introduce babel español (ver CLAUDE.md §7.3.3).
\usetikzlibrary{arrows.meta,calc,angles,quotes}

\tikzset{
  figline/.style={figblue, line width=0.9pt},
  figaux/.style={figgray, dashed, line width=0.6pt},
  % vectores
  figvec/.style={figblue, line width=0.9pt, -{Latex[length=2.2mm,width=1.6mm]}},
  figvecres/.style={figred, line width=1.3pt, -{Latex[length=2.6mm,width=1.9mm]}},
  figvecaux/.style={figgray, line width=0.7pt, -{Latex[length=1.8mm,width=1.3mm]}},
  figaxis/.style={figgray, line width=0.6pt, -{Latex[length=2mm,width=1.4mm]}},
  % cargas puntuales
  figcarga/.style={circle, inner sep=0pt, minimum size=5.4pt, draw=none},
  figpos/.style={figcarga, fill=figred},
  figneg/.style={figcarga, fill=figblue},
  figneutra/.style={figcarga, fill=figgray},
  % etiquetas
  figlbl/.style={font=\small},
  figlblaux/.style={font=\footnotesize, figgray},
}

% ---- pgfplots: estilo base de las gráficas ----
\pgfplotsset{
  emfig/.style={
    width=11cm, height=6.5cm,
    axis lines=left,
    xlabel style={font=\small},
    ylabel style={font=\small},
    tick label style={font=\footnotesize},
    every axis plot/.append style={line width=1pt},
    grid=major,
    grid style={figgray!20},
    legend cell align=left,
    legend style={font=\footnotesize, draw=figgray!40},
  },
}

% (compilar desde el directorio de la clase para que la ruta relativa resuelva).
% ---------------------------------------------------------------------------

% ---- Paleta (alineada con las cajas del preámbulo: keybox/notebox/warnbox) ----
\definecolor{figblue}{HTML}{1F4E79}   % azul principal (líneas, keybox)
\definecolor{figred}{HTML}{C0392B}    % rojo de énfasis (warnbox)
\definecolor{figamber}{HTML}{B8860B}  % ámbar (notebox)
\definecolor{figgray}{HTML}{555555}   % auxiliares, ejes, guías

% ---- CircuiTikz: escala y grosor de los componentes ----
% Guarda \ifdefined: la electrostática (Clases 1-14) no carga `circuitikz`, y
% sin la guarda el \input revienta con `Undefined control sequence \ctikzset`.
% Cuando el curso llegue a circuitos, el bloque se activa solo.
\ifdefined\ctikzset
  \ctikzset{
    bipoles/thickness=1,
    resistors/scale=0.9,
    capacitors/scale=0.9,
  }
\fi

% ---- TikZ: estilos reutilizables ----
% Nota: las puntas se declaran como `-{Latex[...]}` y no como `->`. La forma
% con llaves NO contiene un `>` literal, así que es inmune al `>` activo que
% introduce babel español (ver CLAUDE.md §7.3.3).
\usetikzlibrary{arrows.meta,calc,angles,quotes}

\tikzset{
  figline/.style={figblue, line width=0.9pt},
  figaux/.style={figgray, dashed, line width=0.6pt},
  % vectores
  figvec/.style={figblue, line width=0.9pt, -{Latex[length=2.2mm,width=1.6mm]}},
  figvecres/.style={figred, line width=1.3pt, -{Latex[length=2.6mm,width=1.9mm]}},
  figvecaux/.style={figgray, line width=0.7pt, -{Latex[length=1.8mm,width=1.3mm]}},
  figaxis/.style={figgray, line width=0.6pt, -{Latex[length=2mm,width=1.4mm]}},
  % cargas puntuales
  figcarga/.style={circle, inner sep=0pt, minimum size=5.4pt, draw=none},
  figpos/.style={figcarga, fill=figred},
  figneg/.style={figcarga, fill=figblue},
  figneutra/.style={figcarga, fill=figgray},
  % etiquetas
  figlbl/.style={font=\small},
  figlblaux/.style={font=\footnotesize, figgray},
}

% ---- pgfplots: estilo base de las gráficas ----
\pgfplotsset{
  emfig/.style={
    width=11cm, height=6.5cm,
    axis lines=left,
    xlabel style={font=\small},
    ylabel style={font=\small},
    tick label style={font=\footnotesize},
    every axis plot/.append style={line width=1pt},
    grid=major,
    grid style={figgray!20},
    legend cell align=left,
    legend style={font=\footnotesize, draw=figgray!40},
  },
}


% ---- Title ----
\title{\textbf{Clase 28: Circuitos RLC con mutuas y ecuaciones de Maxwell} \\[0.3em]
        \large{Resonancia, convención de los puntos, transformador y corriente de desplazamiento}}
\author{Transcripto y expandido --- Curso Electromagnetismo (OpenFING)}
\date{}

\begin{document}

\maketitle
\thispagestyle{empty}
\newpage

\tableofcontents
\newpage

\section{Resonancia}

\subsection{La idea: la hamaca}

La \textbf{resonancia} es una propiedad general de cualquier sistema oscilante
---eléctrico o mecánico---: con un forzado a ciertas frecuencias particulares, el
oscilador oscila con amplitud \textbf{muy grande}.

El ejemplo de siempre del docente: una hamaca. (No es un oscilador lineal para
ángulos grandes, pero se piensa con oscilaciones chicas.) Sin forzado y con poco
rozamiento tiene una \textbf{frecuencia natural}: digamos, va y vuelve una vez por
segundo. Empujarla cuatro veces por segundo no es óptimo ---a veces se empuja a
contramano---; empujarla una vez cada cuatro segundos desaprovecha oportunidades.
Lo ideal es empujarla \textbf{cada vez que vuelve}: con una frecuencia igual o muy
parecida a la natural.

\begin{notebox}
La analogía es cualitativa: la hamaca es no lineal y el empujón es periódico
pero no sinusoidal (más bien un golpe). Lo que se conserva es la idea.
\end{notebox}

El resultado que se va a ver: la frecuencia de máxima amplitud coincide con la
natural \textbf{sólo sin rozamiento} (sin resistencia); con resistencia son
parecidas, no iguales.

\subsection{Frecuencia de resonancia del RLC serie}

El RLC serie con $V_0\cos\omega t$:

\[
I(t) = \frac{V_0}{|\hat Z|}\cos(\omega t-\theta),
\qquad
|\hat Z| = \sqrt{R^2 + \left(\omega L - \frac{1}{\omega C}\right)^2},
\qquad
\theta = \arctan\frac{\omega L - 1/\omega C}{R}.
\]

La amplitud de la corriente, $I_0(\omega) = V_0/|\hat Z|$, es máxima cuando
$|\hat Z|$ es mínimo. A simple vista: $|\hat Z|^2$ es una suma de dos
cuadrados, $R^2$ no depende de $\omega$, y el otro se puede anular. Por lo
tanto $|\hat Z|\ge R$, con igualdad cuando \textbf{la reactancia se anula}:

\[
\omega L = \frac{1}{\omega C}
\quad\Longrightarrow\quad
\boxed{\ \omega_0 = \frac{1}{\sqrt{LC}}\ },
\qquad
I_0(\omega_0) = \frac{V_0}{R}.
\]

La \textbf{curva de resonancia} $I_0(\omega)$ arranca en cero (para $\omega\to0$ el
condensador hace divergir $|\hat Z|$), tiene un máximo $V_0/R$ en $\omega_0$ y
vuelve a tender a cero para $\omega\to\infty$ (diverge el término de la
bobina).

\subsection{Frecuencia natural y frecuencia de resonancia}

Sin forzado, el circuito oscila a la \textbf{frecuencia natural} calculada dos
clases atrás,

\[
\omega_0' = \sqrt{\frac{1}{LC}-\frac{R^2}{4L^2}} .
\]

Cuando $R\to0$, $\omega_0'\to\omega_0$: con resistencias pequeñas, el forzado
óptimo está en la frecuencia natural. Con resistencia no nula, la \textbf{frecuencia
de resonancia} $\omega_0$ y la natural no coinciden.

\begin{notebox}
Un estudiante observa que la que se mueve con $R$ es $\omega_0'$, no
$\omega_0$, que no depende de $R$. El docente lo acepta: lo que quería decir es
que se acercan.
\end{notebox}

Además, al bajar $R$ el máximo $V_0/R$ crece: el pico se vuelve cada vez más
alto y angosto. En el límite $R\to0$ sería infinito, situación no física; pero
con muy poca resistencia la resonancia puede ser muy grande.

\begin{notebox}[Abuso de lenguaje estándar]
En todo lo que sigue se dice «frecuencia» por
\textbf{frecuencia angular}, que difiere en $2\pi$.
\end{notebox}

\subsection{Ancho de la resonancia y factor de calidad}

Hay que definir qué se llama ancho; la convención es tomar las frecuencias en
que la amplitud cae al máximo \textbf{dividido por $\sqrt2$} (es arbitrario, pero las
cuentas quedan cómodas):

\[
\frac{I_0(\omega)}{I_0(\omega_0)} = \frac{1}{\sqrt2}
\quad\Longleftrightarrow\quad
|\hat Z(\omega)| = \sqrt2\,R
\quad\Longleftrightarrow\quad
\left|\omega L - \frac{1}{\omega C}\right| = R .
\]

\begin{notebox}
El docente pone primero el $\sqrt2$ del lado equivocado y lo corrigen.
\end{notebox}

Se calcula sólo para \textbf{$R$ pequeño}, donde el ancho es chico:
$\omega = \omega_0 + \Delta\omega$ con $\Delta\omega\ll\omega_0$. Con
$(1+x)^{-1}\approx 1-x$ («dividir por 1,01 es como multiplicar por 0,99»):

\[
\frac{1}{C(\omega_0+\Delta\omega)} \approx \frac{1}{C\omega_0}\left(1-\frac{\Delta\omega}{\omega_0}\right).
\]

El término de orden cero se cancela por la condición de resonancia, y como
$1/(C\omega_0^2) = L$:

\[
L\,\Delta\omega + L\,\Delta\omega = R
\quad\Longrightarrow\quad
\boxed{\ \Delta\omega = \frac{R}{2L}\ }
\]

$\Delta\omega$ es el ancho hacia cada lado; cerca del pico, con $R$ chico, la
curva es simétrica. Tiende a cero con $R$.

Se define el \textbf{factor de calidad}

\begin{keybox}
\[
\boxed{\ \mathcal{Q} = \frac{\omega_0 L}{R} = \frac{\omega_0}{2\Delta\omega}\ }
\]
\end{keybox}

(con la letra Q «porque se me acabaron las letras»; no es la carga). Cuanto
mayor, más aguda la resonancia.

\begin{notebox}[Por qué «calidad»]
Una radio: la antena recibe una señal débil, que actúa
como un generador de forzado pequeño. Las radios reales tienen amplificador,
pero aun sin él, con un $\mathcal{Q}$ alto y bien sintonizada, ese forzado
pequeño produce oscilaciones grandes. Un estudiante resume: es más selectivo
y además amplifica mucho en resonancia. El costo es que exige una
sintonización más fina. Las radios viejas sintonizaban con un condensador o
una resistencia variable; hoy es todo electrónico, porque se miniaturiza
mejor, aunque la física de fondo es la misma.
\end{notebox}

\subsection{La resonancia en la fase}

La resonancia también se ve en la fase, $\theta = \arctan\dfrac{\omega L - 1/\omega C}{R}$:

\begin{center}
\renewcommand{\arraystretch}{1.35}
\small
\begin{tabularx}{\textwidth}{l>{\raggedright\arraybackslash}X>{\raggedright\arraybackslash}X}
\toprule
\textbf{$\omega$} & \textbf{Argumento} & \textbf{$\theta$} \\
\midrule
$\to0$ & $\to-\infty$ & $\to-\pi/2$ \\
$\omega_0$ & $0$ & $0$ \\
$\to\infty$ & $\to+\infty$ & $\to+\pi/2$ \\
\bottomrule
\end{tabularx}
\end{center}

La fase pasa de $-\pi/2$ a $\pi/2$ en una región de ancho del orden de
$\Delta\omega$. Al bajar $R$ el pasaje se vuelve más vertical, y en el límite
$R\to0$ es una discontinuidad.

\begin{figure}[H]
  \centering
  % Clase 28 — curvas de resonancia del RLC serie para tres resistencias, en
% unidades de omega0 = 1/sqrt(LC). Con x = omega/omega0 y Q = omega0 L / R:
%   amplitud  I0/(V0/R) = 1/sqrt(1 + Q^2 (x - 1/x)^2)   ... (normalizada al
%             máximo de cada curva; el pico absoluto V0/R crece al bajar R)
%   fase      theta = arctan(Q (x - 1/x))
\begin{tikzpicture}
  \begin{axis}[
    emfig, width=6.9cm, height=5.4cm,
    xmin=0.3, xmax=2, ymin=0, ymax=1.1,
    xtick={0.5,1,1.5,2}, ytick={0,0.5,0.707,1},
    yticklabels={$0$,$0{,}5$,$1/\sqrt2$,$1$},
    xlabel={$\omega/\omega_0$}, ylabel={$I_0/I_0(\omega_0)$},
    title={\small (a) amplitud},
    legend style={at={(1.1,-0.42)}, anchor=north, legend columns=3,
      /tikz/every even column/.append style={column sep=1em}},
  ]
    \addplot[figblue!40, domain=0.3:2, samples=200] {1/sqrt(1+4*(x-1/x)^2)};
    \addlegendentry{$\mathcal{Q} = 2$}
    \addplot[figblue!75, domain=0.3:2, samples=200] {1/sqrt(1+25*(x-1/x)^2)};
    \addlegendentry{$\mathcal{Q} = 5$}
    \addplot[figblue, domain=0.3:2, samples=400] {1/sqrt(1+400*(x-1/x)^2)};
    \addlegendentry{$\mathcal{Q} = 20$}
  \end{axis}
  \begin{axis}[
    emfig, width=6.9cm, height=5.4cm, xshift=7.6cm,
    xmin=0.3, xmax=2, ymin=-1.7, ymax=1.7,
    xtick={0.5,1,1.5,2}, ytick={-1.5708,0,1.5708},
    yticklabels={$-\pi/2$,$0$,$\pi/2$},
    xlabel={$\omega/\omega_0$}, ylabel={$\theta$},
    title={\small (b) fase},
  ]
    \addplot[figblue!40, domain=0.3:2, samples=200] {rad(atan(2*(x-1/x)))};
    \addplot[figblue!75, domain=0.3:2, samples=200] {rad(atan(5*(x-1/x)))};
    \addplot[figblue, domain=0.3:2, samples=400] {rad(atan(20*(x-1/x)))};
  \end{axis}
\end{tikzpicture}

  \caption{Curvas de resonancia del RLC serie para tres factores de calidad: (a) amplitud, normalizada a su máximo $V_0/R$, que crece al bajar $R$; (b) fase, que pasa de $-\pi/2$ a $\pi/2$ cada vez más bruscamente.}
\end{figure}

\subsection{Varios osciladores acoplados}

El fenómeno es general. Con $n$ osciladores lineales acoplados y un generador
sinusoidal, \textbf{para resistencias pequeñas}:

\begin{enumerate}
  \item hay \textbf{$n$ frecuencias de resonancia} ($n$ picos);
  \item cada una es cercana a una de las \textbf{$n$ frecuencias naturales} del sistema
  sin forzar, y en el límite de resistencias nulas coinciden;
  \item en resonancia el desfasaje se anula (estrictamente, en el límite $R\to0$) y
  el \textbf{factor de potencia tiende a 1}; con resistencias pequeñas, «factor de
  potencia 1» y «módulo de la impedancia mínimo» son equivalentes.
\end{enumerate}

\begin{notebox}[Por qué resistencias pequeñas]
Con resistencias grandes el sistema sin
forzar ni siquiera oscila: decae sin oscilar, y la noción de frecuencia
natural pierde sentido. Cuanto mayor la resistencia, más chatos los picos,
que pueden volverse invisibles.
\end{notebox}

Se puede probar en general; el docente no lo hace.

\section{Inductancias mutuas en circuitos}

\subsection{La convención de los puntos}

El signo de una mutua depende de cómo están enrolladas las bobinas: dos
bobinados en el mismo solenoide enrollados para el mismo lado dan $M>0$; al
revés, $M<0$. En un diagrama de circuito eso no se ve. La convención:

\begin{notebox}
\textbf{Se da $M$ siempre positiva y el signo se indica en el diagrama con puntos}
en un extremo de cada bobina.
\end{notebox}

\textbf{Regla.} Si todas las corrientes (elegidas convencionalmente) \textbf{entran por
los puntos}, los términos de mutua llevan \textbf{signo $+$}, el mismo que las
autoinductancias. Cada cambio de un punto al otro extremo introduce un signo
delante de $j\omega M$; dos cambios se compensan. Si las dos corrientes entran
por el extremo sin punto, también es $+$.

\begin{notebox}
Con varias mutuas no siempre se puede elegir que todas las corrientes entren
por los puntos: al acomodar una, otra puede quedar al revés.
\end{notebox}

\subsection{Ejemplo: dos ramas acopladas}

Un generador alimenta dos ramas en paralelo: una con $\hat Z_1$ y $L_1$, otra con
$\hat Z_2$ y $L_2$, acopladas por $M$, con los dos puntos del mismo lado y las
dos corrientes entrando por ellos. Malla de arriba (con el generador) y malla
de abajo:

\[
V_0 = \hat Z_1\hat I_1 + j\omega L_1\hat I_1 + j\omega M\,\hat I_2,
\]

\[
0 = \hat Z_2\hat I_2 + j\omega L_2\hat I_2 + j\omega M\,\hat I_1 - (\text{rama de arriba}).
\]

\begin{notebox}
En la segunda malla el docente escribe primero $V_0$ del lado izquierdo y se
corrige: «mezclé dos cosas». La malla de abajo no contiene el generador. El
punto importante es el signo de los términos con $M$, que es $+$ en las dos
porque ambas corrientes entran por los puntos.
\end{notebox}

Es un sistema lineal que se puede resolver; el docente no lo resuelve.

\section{El transformador}

\subsection{Transformador real}

\begin{figure}[H]
  \centering
  % Clase 28 — transformador: circuito primario (generador, R1, L1) y secundario
% (L2, R2, carga Z), acoplados por M. Los puntos indican que, con las dos
% corrientes entrando por ellos, los términos de mutua llevan signo +.
\begin{tikzpicture}[scale=1]
  \begin{scope}[line width=0.9pt, color=figblue]
    % primario
    \draw (0,0) to[sV] (0,2.6);
    \draw (0,2.6) to[R] (2.2,2.6) -- (2.8,2.6);
    \draw (2.8,2.6) to[L] (2.8,0);
    \draw (2.8,0) -- (0,0);
    % secundario
    \draw (4.0,2.6) to[L] (4.0,0);
    \draw (4.0,2.6) -- (4.6,2.6) to[R] (6.6,2.6) -- (6.8,2.6);
    \draw (6.8,2.6) to[generic] (6.8,0);
    \draw (6.8,0) -- (4.0,0);
  \end{scope}
  % puntos
  \fill[figblue] (2.45,2.25) circle (2pt);
  \fill[figblue] (4.35,2.25) circle (2pt);
  % rotulos
  \node[figlbl, anchor=east] at (-0.55,1.3) {$V_0\cos\omega t$};
  \node[figlbl, figblue] at (1.1,3.05) {$R_1$};
  \node[figlbl, figblue, anchor=east] at (2.35,1.1) {$L_1$};
  \node[figlbl, figblue, anchor=west] at (4.45,1.1) {$L_2$};
  \node[figlbl, figblue] at (5.6,3.05) {$R_2$};
  \node[figlbl, figblue, anchor=west] at (7.25,1.3) {$\hat Z$};
  \draw[figgray, line width=0.6pt, {Latex[length=1.6mm]}-{Latex[length=1.6mm]}] (2.95,1.6) -- (3.85,1.6);
  \node[figlbl] at (3.4,1.9) {$M$};
  % corrientes entrando por los puntos
  \draw[figvecres] (1.9,2.95) -- (2.55,2.95);
  \node[figlbl, text=figred] at (2.2,3.3) {$\hat I_1$};
  \draw[figvecres] (5.0,2.95) -- (4.35,2.95);
  \node[figlbl, text=figred] at (4.7,3.3) {$\hat I_2$};
  \node[figlblaux, anchor=north] at (1.4,-0.2) {primario};
  \node[figlblaux, anchor=north] at (5.4,-0.2) {secundario};
\end{tikzpicture}

  \caption{Transformador: primario con el generador, $R_1$ y $L_1$; secundario con $L_2$, $R_2$ y la carga $\hat Z$. Con las dos corrientes entrando por los puntos, los términos de mutua llevan signo $+$.}
\end{figure}

Un \textbf{circuito primario} con el generador $V_0\cos\omega t$, una resistencia
$R_1$ y una bobina $L_1$, y un \textbf{circuito secundario}, sin generador, con una
bobina $L_2$, una resistencia $R_2$ y la carga $\hat Z$ que se quiere
alimentar. $R_1$ y $R_2$ representan las resistencias internas, inevitables, de
los bobinados. Con las corrientes entrando por los puntos:

\[
V_0 = (R_1 + j\omega L_1)\,\hat I_1 + j\omega M\,\hat I_2,
\]

\[
0 = j\omega M\,\hat I_1 + (R_2 + j\omega L_2 + \hat Z)\,\hat I_2 .
\]

Resolviendo el sistema $2\times2$, con
$D = (R_1 + j\omega L_1)(\hat Z + j\omega L_2 + R_2) + \omega^2M^2$:

\begin{keybox}
\[
\boxed{\ \hat I_1 = \frac{\hat Z + j\omega L_2 + R_2}{D}\,V_0,
\qquad
\hat I_2 = -\frac{j\omega M}{D}\,V_0\ }
\]
\end{keybox}

\begin{itemize}
  \item Si $M\to0$ (bobinas desacopladas), en el secundario no pasa nada: la
  corriente ahí es puramente inducida.
  \item Las dos corrientes son proporcionales a $V_0$.
  \item Es la solución exacta de un \textbf{transformador real}: no se supuso acoplamiento
  máximo ($M$ puede ser menor que $\sqrt{L_1L_2}$) y hay resistencias.
\end{itemize}

No hace lo que se quiere de un transformador: los cocientes de corrientes y de
voltajes \textbf{dependen de la frecuencia}.

\subsection{Transformador ideal}

Un transformador ideal cambia voltaje y corriente por una constante
\textbf{independiente de $\omega$}: multiplica el voltaje por un factor $A$ y la
corriente por el inverso.

\begin{notebox}
En el VTT el docente dice que $\hat I_2/\hat I_1$ y $\hat V_Z/V_0$ son «la misma
constante» $A$ (con un signo menos en la corriente), y enseguida que la
corriente cambia «de manera inversa». Por el resultado de abajo y el del
ejercicio, lo consistente es: voltaje por $A$, corriente por $1/A$.
\end{notebox}

Del sistema:

\[
-\frac{\hat I_1}{\hat I_2} = \frac{\hat Z + j\omega L_2 + R_2}{j\omega M},
\]
que se vuelve independiente de $\omega$ si

\begin{keybox}
\[
\boxed{\ \omega L_2 \gg |\hat Z + R_2|\ }
\quad\Longrightarrow\quad
-\frac{\hat I_1}{\hat I_2}\approx\frac{L_2}{M}.
\]
\end{keybox}

Es la primera condición de transformador ideal: resistencias pequeñas y una
carga que no le exija demasiado.

\textbf{Ejercicio propuesto.} Analizar el voltaje de manera similar, en el caso de
\textbf{acoplamiento máximo}, $M^2 = L_1L_2$ (por ejemplo, dos bobinados en el mismo
solenoide), y probar que

\[
A = \frac{N_2}{N_1},
\]
con $N_1$ y $N_2$ los números de espiras.

\begin{notebox}
Con $L\propto N^2$ y $M = \sqrt{L_1L_2}$, el cociente de corrientes de arriba
da $L_2/M = N_2/N_1$, coherente con que la corriente cambie por el inverso de
$A$.
\end{notebox}

\section{Ecuaciones de Maxwell}

Capítulo 11, el último del curso.

\subsection{La inconsistencia de la ley de Ampère}

Un cable con corriente variable $I(t)$ que llega a un condensador. Se aplica
Ampère a una curva $C$ alrededor del cable, con dos superficies del mismo borde:
$S_1$, atravesada por el cable, y $S_2$, que «da la vuelta» pasando entre las
placas.

\[
\oint_C \vec H\cdot d\vec l = \int_S \vec J\cdot\hat n\,dS =
\begin{cases} I & \text{con } S_1\\ 0 & \text{con } S_2 \end{cases}
\]

\begin{figure}[H]
  \centering
  % Clase 28 — la inconsistencia de Ampère: un cable con corriente I(t) llega a
% un condensador. Con la misma curva C, la superficie S1 la cruza el cable
% (da I) y la S2 pasa entre las placas (daría 0). La diferencia es la carga que
% se acumula en la placa.
\begin{tikzpicture}[scale=1]

  % cable y placas
  \draw[figblue, line width=1.6pt] (-3.0,0) -- (1.2,0);
  \draw[figblue, line width=1.6pt] (2.4,0) -- (4.2,0);
  \draw[figblue, line width=2.2pt] (1.2,-1.3) -- (1.2,1.3);
  \draw[figblue, line width=2.2pt] (2.4,-1.3) -- (2.4,1.3);
  \node[figlblaux, text=figred, anchor=south] at (1.0,1.95) {$+Q(t)$};
  \node[figlblaux, text=figblue, anchor=south] at (2.6,1.95) {$-Q(t)$};
  \draw[figvecres] (-2.6,0.25) -- (-1.8,0.25);
  \node[figlbl, text=figred] at (-2.2,0.6) {$I(t)$};

  % curva C (vista de costado: una elipse alrededor del cable)
  \draw[figred, line width=1pt] (-0.6,0) ellipse (0.18 and 0.95);
  \node[figlbl, text=figred] at (-0.6,-1.25) {$C$};

  % S1: disco plano en la curva
  \fill[figamber!25, opacity=0.8] (-0.6,0) ellipse (0.18 and 0.95);
  \node[figlbl, text=figamber] at (-1.15,0.95) {$S_1$};

  % S2: bolsa que rodea la placa +Q y se cierra entre las placas
  \draw[figamber, line width=0.9pt, dashed]
    (-0.6,0.95) .. controls (0.2,1.75) and (1.8,1.9) .. (1.8,0.9)
               -- (1.8,-0.9)
               .. controls (1.8,-1.9) and (0.2,-1.75) .. (-0.6,-0.95);
  \node[figlbl, text=figamber, fill=white, inner sep=1pt] at (1.8,0.45) {$S_2$};

\end{tikzpicture}

  \caption{La inconsistencia de la ley de Ampère: con la misma curva $C$, el cable atraviesa $S_1$, mientras que $S_2$ rodea la placa y pasa entre las placas, donde no hay corriente de conducción.}
\end{figure}

Tendrían que dar lo mismo y no dan. No es extraño: Ampère se dedujo para
corrientes estacionarias.

La diferencia entre los dos cálculos es el flujo de $\vec J$ por la superficie
cerrada formada por $S_1$ y $S_2$, que es la corriente que entra en ella, igual
a la variación de la carga encerrada:

\[
\int_{S_1}\vec J\cdot\hat n\,dS - \int_{S_2}\vec J\cdot\hat n\,dS
= \frac{d}{dt}\int_V\rho\,dV .
\]

El motivo de la inconsistencia: \textbf{la carga en el condensador está variando.}

\subsection{Corregirla con la conservación de la carga}

En forma diferencial: la divergencia de un rotacional es cero, así que

\[
\curl\vec H = \vec J
\quad\Longrightarrow\quad
\div\vec J = 0 :
\]
la ecuación \textbf{sólo es consistente si $\div\vec J = 0$}, es decir, en el caso
estacionario.

Hay que reemplazar $\vec J$ por un nuevo campo $\vec J_{\text{nuevo}}$ tal que

\begin{enumerate}
  \item $\vec J_{\text{nuevo}} = \vec J$ en el caso estacionario;
  \item $\div\vec J_{\text{nuevo}} = 0$ \textbf{siempre}.
\end{enumerate}

No hay una única solución; se busca la más simple, que es la que encontró
Maxwell y que además resultó la correcta, aunque «podría haber sido más
complicada».

La pista es la \textbf{conservación de la carga}. Lo que se sabe de $\div\vec J$ en
el caso general es la ecuación de continuidad, y lo que se sabe de $\rho$ en
términos de una divergencia es la ley de Gauss:

\[
\pdv{\rho}{t} + \div\vec J = 0,
\qquad
\div\vec D = \rho
\quad\Longrightarrow\quad
\div\left(\vec J + \pdv{\vec D}{t}\right) = 0 .
\]

$\vec J + \partial\vec D/\partial t$ cumple las dos condiciones: en el caso
estacionario es $\vec J$, y su divergencia es siempre nula (se usaron sólo
ecuaciones exactas).

\subsection{Corriente de desplazamiento}

Maxwell propone modificar la ley de Ampère:

\begin{keybox}
\[
\boxed{\ \curl\vec H = \vec J + \pdv{\vec D}{t}\ }
\]
\end{keybox}

\textbf{Ley de Ampère modificada por Maxwell.} (El docente dice «divergencia» en voz
alta y escribe rotacional; lo corrige.) A $\partial\vec D/\partial t$ se lo llama
\textbf{densidad de corriente de desplazamiento}: en Física III se hablaba de
corriente de desplazamiento sin establecer el vínculo con el \textbf{vector
desplazamiento} $\vec D$, que es de donde viene el nombre. La \textbf{corriente de
desplazamiento} es

\[
I_d = \int_S \pdv{\vec D}{t}\cdot\hat n\,dS .
\]

\begin{notebox}
A la pregunta de si Maxwell «se llevó todo el crédito»: «a Ampère le seguimos
diciendo Ampère».
\end{notebox}

\subsection{Las cuatro ecuaciones}

\begin{center}
\renewcommand{\arraystretch}{1.35}
\small
\begin{tabularx}{\textwidth}{l>{\raggedright\arraybackslash}X>{\raggedright\arraybackslash}X}
\toprule
\textbf{Ley} & \textbf{En un medio} & \textbf{En el vacío} \\
\midrule
Gauss & $\div\vec D = \rho$ & $\div\vec E = \rho/\varepsilon_0$ \\
No hay monopolos & $\div\vec B = 0$ & $\div\vec B = 0$ \\
Faraday & $\curl\vec E = -\pdv{\vec B}{t}$ & $\curl\vec E = -\pdv{\vec B}{t}$ \\
Ampère--Maxwell & $\curl\vec H = \vec J + \pdv{\vec D}{t}$ & $\curl\vec B = \mu_0\vec J + \mu_0\varepsilon_0\pdv{\vec E}{t}$ \\
\bottomrule
\end{tabularx}
\end{center}

«Está lleno de remeras que tienen esto.» En un medio hay que agregar siempre
las relaciones constitutivas entre $\vec D$ y $\vec E$ y entre $\vec H$ y
$\vec B$.

\subsection{¿Siguen siendo correctas?}

Maxwell hizo una jugada arriesgada: buscó «al tuntún» la corrección más simple,
y era la correcta. Hoy se piensa que las ecuaciones están bien, con dos
modificaciones del siglo XX:

\begin{enumerate}
  \item \textbf{Efectos cuánticos.} A escalas muy chicas hay que usar el equivalente
  cuántico de estas ecuaciones: la \textbf{electrodinámica cuántica}, que lograron
  formular de manera coherente Feynman, Schwinger y Tomonaga (premio Nobel).
  No cambia las ecuaciones, sino cómo se usan a nivel microscópico; para
  objetos macroscópicos el límite es esto.
  \item \textbf{El modelo estándar.} A distancias del orden de un centésimo del radio del
  protón, el electromagnetismo se entrevera con las interacciones nucleares
  débiles: la \textbf{unificación electrodébil}.
\end{enumerate}

Fuera de esas escalas, son leyes fundamentales del electromagnetismo, y con
enormes consecuencias prácticas: tienen \textbf{soluciones en forma de ondas}. Así
Maxwell entendió la naturaleza electromagnética de la luz y pudo calcular su
velocidad a partir de propiedades eléctricas y magnéticas.

\noindent\emph{La próxima clase: las ondas electromagnéticas.}

\end{document}
